Autoregressive models generate sequences monotonically, where any sampled token, even if erroneous or sub-optimal, becomes a permanent condition for all subsequent steps.
This structural limitation means that autoregressive models cannot revisit or revise earlier decisions, i.e., a capability essential for complex generation tasks where exploration and correction are necessary.
To this end, we propose \textbf{\modelname}, a \textbf{N}on-\textbf{M}onotonic \textbf{A}utoreg\textbf{R}essive \textbf{S}equence modeling framework that enables models to generate, evaluate, and revise tokens within a single forward pass, effectively allowing exploration before commitment.
We operationalize this framework through a learned erase token \bk that retracts the previous token, enabling on-the-fly revision within standard autoregressive decoding.
To train the model, we introduce a sequence augmentation method that constructs error-correction trajectories from model-generated deviations paired with ground-truth references.
We then propose masked supervised fine-tuning (mSFT) that exposes the model to errors as context for learning when to revise, without optimizing their likelihood.
Finally, we refine the model with group relative policy optimization (GRPO) that incentivizes judicious \bk usage, i.e., rewarding effective corrections while penalizing unsuccessful ones. 
We conduct comprehensive theoretical and empirical analysis to validate the effectiveness of \modelname, demonstrating a robust foundation for non-monotonic sequence modeling.